\newpage\section{\tkzClass{matrix} 类}

变量 \tkzVar{matrix}{M} 保存一个用于存储矩阵的表。它是可选的，你可以自由选择变量名。不过，为清晰与一致起见，推荐使用 \code{M}。如果你使用自定义变量（例如 Matrices），则必须手动初始化它。\\
若使用 \code{init\_elements()} 函数，它会重新初始化 \code{M} 表。

\vspace{1em}

\code{matrix} 类目前是实验性的，其属性名与方法名尚未最终确定，说明该类仍在演进中。它已与其他类（例如 \code{point} 类）建立了某些联系。此外还加入了一个新属性 \code{mtx}，把一个列矩阵与点关联起来，其中的元素对应于该点在原基下的坐标。类似地，\code{vector} 类也加入了一个属性，其中 \code{mtx} 表示由构成该向量的两个附标组成的列矩阵。

创建这个 \code{matrix} 类是为了避免依赖外部库，并已适配平面变换。它允许你使用复数。

\lefthand\ 要显示矩阵，需要加载 \tkzNamePack{amsmath} 宏包。

{\color{red}\lefthand\ } 虽然有些方法适用于任意大小的矩阵，但大多数方法仅适用于 2 阶和 3 阶方阵。


\subsection{矩阵的创建}
\label{sub:matrix_creation}

创建矩阵有多种方式。下面我们逐一查看不同情形。

第一种是使用数组的数组，即一个表，其中每个元素又是另一个表。例如，可用下面的代码创建一个 N 行 M 列的零矩阵：


\begin{itemize}

\item  第一种方法见 [\ref{ssub:method_new}]。这个函数最重要，因为它负责创建对象。其他函数创建特定对象时总是使用这个函数。

\tkzClass{matrix} 类表示由四个值定义的 $2\times2$ 矩阵。

\begin{mybox}
\begin{verbatim}
M = matrix:new(1, 0, 0, 1) -- identity matrix
\end{verbatim}
\end{mybox}

\texttt{简写形式：}

还有一种更简洁的形式：

\begin{mybox}
\begin{verbatim}
M = matrix(1, 0, 0, 1)
\end{verbatim}
\end{mybox}


\begin{mybox}
   | M.new = matrix({ { a, b }, { c, d } }) | \\
   其中 a、b、c、d 为实数或复数。
\end{mybox}

\begin{minipage}{.3\textwidth}
\directlua{
  init_elements()
  local a, b, c, d = 1, 2, 3, 4
  M.new = matrix({ { a, b }, { c, d } })
  tex.print('M = ') M.new:print()}
\end{minipage}
\begin{minipage}{.6\textwidth}
\begin{tkzexample}[code only]
\directlua{
  init_elements()
  local a, b, c, d = 1, 2, 3, 4
  M.new = matrix({ { a, b }, { c, d } })
  tex.print('M = ') M.new:print()}
\end{tkzexample}
\end{minipage}


\item 使用函数 \code{create}，可得到一个系数全为零、行列数可自选的矩阵。[\ref{ssub:function_matrix_create_n_m}]

  \begin{mybox}
  |M.cr = matrix.create(4,5)|
  \end{mybox}

\begin{tkzexample}[latex=.5\textwidth]
\directlua{
  init_elements()
  M.cr = matrix.create(4, 5)
  tex.print('M = ') M.cr:print()}
\end{tkzexample}


\item  大小为 {$\displaystyle n$} 的单位矩阵是一个 {$\displaystyle n\times n$} 方阵，主对角线上为 1、其余位置为 0。见 [\ref{ssub:method_identity}]

\begin{mybox}
|M.I = matrix.identity(3)|
\end{mybox}


\begin{tkzexample}[latex=.5\textwidth]
\directlua{
  init_elements()
  M.I = matrix.identity(3)
  tex.print('$I_3 = $') M.I:print()}
\end{tkzexample}


\item  也可以用下面的方式得到方阵：[\ref{ssub:method_square}]
  \begin{mybox}
  |M.sq = matrix.square (2,a,b,c,d)|
  \end{mybox}

\begin{tkzexample}[latex=.5\textwidth]
\directlua{
  init_elements()
   local a, b, c, d = 1, 2, 3, 4
  M.sq = matrix.square(2, a, b, c, d)
  tex.print('M = ') M.sq:print()}
\end{tkzexample}


\item  列向量的情形：[\ref{ssub:method_vector}]

  \begin{mybox}
  | M.V = matrix.vector(1, 2, 3)|
  也可以用 |M.V = matrix.column(1, 2, 3)|
  \end{mybox}


\begin{tkzexample}[latex=.5\textwidth]
\directlua{
  init_elements()
  M.V = matrix.vector(1, 2, 3)
  tex.print('V = ') M.V:print()}
\end{tkzexample}


\item  行向量的情形：[\ref{ssub:function_row_vector}]

\begin{mybox}
  | M.V = matrix.row_vector(1, 2, 3)|
\end{mybox}

\begin{tkzexample}[latex=.5\textwidth]
\directlua{
  init_elements()
  M.V = matrix.row_vector(1, 2, 3)
  tex.print('V = ') M.V:print()}
\end{tkzexample}


\item  与点关联的矩阵

  |M.p = matrix({ { p.re }, { p.im } })|

\item  与向量关联的矩阵

 它是一个列矩阵，由定义该向量的两个点的附标组成。

| local M.v = matrix{ { za }, { zb } }|

\begin{tkzexample}[latex=.5\textwidth]
\directlua{
z.A = point(1, 2)
z.B = point(3, 4)
V.u = vector(z.A, z.B)
V.u.mtx:print()}
\end{tkzexample}



\item  齐次变换矩阵 [\ref{ssub:method_htm}]

 目的是生成一个带齐次坐标的矩阵，能够通过旋转、平移和缩放来变换坐标系。为此需要同时指定旋转角、新原点的坐标以及缩放因子。

\begin{tkzexample}[latex=.35\textwidth]
\directlua{
  init_elements()
  M.h = matrix.htm(math.pi / 3, 1, 2, 2, 1)
  tex.print('H = ') M.h:print()}
\end{tkzexample}

\end{itemize}

\subsection{方法 \tkzMeth{matrix}{print()}}
\label{sub:display_a_matrix_method_print}


这个方法（见 \ref{ssub:method_print}）用于查看结果，下面说明如何使用它。它可用于实矩阵或复矩阵、方阵或非方阵。有几个选项可用于格式化结果。要使用 “print” 方法需要加载 \tkzNamePack{amsmath} 宏包。若没有该宏包，可以用 |print_array (M)| 不加格式地显示矩阵内容。

\begin{tkzexample}[latex=.5\textwidth]
\directlua{
  init_elements()
  M.new = matrix { { 1, -1}, { 2, 0 } }
  M.new:print()}
\end{tkzexample}

\subsection{矩阵的属性}
\label{sub:attibutes_of_a_matrix}

\begin{center}
  \bgroup
  \catcode`_=12
  \small
  \captionof{table}{矩阵的属性。}\label{matrix:attributes}
  \begin{tabular}{ll}
  \toprule
  \textbf{属性}     & \texttt{参考} \\
  \midrule
  \tkzAttr{matrix}{set}     & [\ref{sub:attribute_set}]\\
  \tkzAttr{matrix}{rows}    & [\ref{ssub:attributes_matrix_rows_and_cols}] \\
  \tkzAttr{matrix}{cols}    & [\ref{ssub:attributes_matrix_rows_and_cols}] \\
  \tkzAttr{matrix}{type}    &   \\
  \tkzAttr{matrix}{det}     & [\ref{ssub:attributes_matrix_det}]\\
  \bottomrule %
  \end{tabular}
  \egroup
\end{center}



\subsubsection{属性 \tkzAttr{matrix}{type}}
\label{ssub:attribute_iattr_matrix_type}
\begin{mybox}
  |M.new = matrix{ { 1, 1}, { 0, 2} } |
  |A = { { 1, 1 }, { 0, 2 } } |
\end{mybox}

\code{M} 是一个矩阵（因而也是一个表），而 A 是一个表。因此 \code{M.type} 给出 \code{'matrix'}，而 \code{A.type = nil}。\code{type(A) 或 type(M) = table}。

\subsubsection{属性 \tkzAttr{matrix}{set}}
\label{sub:attribute_set}
像 |{{1,2},{2,-1}}| 这样的简单数组通常被视为 \code{matrix}。在 \tkzNamePack{tkz-elements} 中，我们把由
|matrix({ { 1, 1 }, { 0, 2 } })|
定义的 |M.new| 视为矩阵，而把 |M.new.set| 视为数组（|M.new.set = { { 1, 1 }, {0, 2 } }|）。

你可以访问矩阵的某个特定元素，例如：|M.new.set[2][1]| 给出 \tkzUseLua{M.new.set[2][2]}。

|\tkzUseLua{M.new.set[2][1]}| 是显示 $2$ 的表达式。

\subsubsection{属性 \tkzAttr{matrix}{rows} 与 \tkzAttr{matrix}{cols}}
\label{ssub:attributes_matrix_rows_and_cols}

行数用 |M.n.rows| 访问，列数用 |M.n.cols| 访问，示例如下：

\vspace{.5em}
\begin{minipage}{.5\textwidth}
\begin{verbatim}
\directlua{
 init_elements()
 M.n = matrix({ { 1, 2, 3 }, { 4, 5, 6 } })
 M.n:print()
 tex.print("Rows:  "..M.n.rows)
 tex.print("Cols:  "..M.n.cols)}
\end{verbatim}
\end{minipage}
\begin{minipage}{.5\textwidth}
\directlua{
 init_elements()
 M.n = matrix({ { 1, 2, 3 }, { 4, 5, 6 } })
 M.n:print()
 tex.print("Rows:  "..M.n.rows)
 tex.print("Cols:  "..M.n.cols)}
\end{minipage}

\subsubsection{属性 \tkzAttr{matrix}{det} }
\label{ssub:attributes_matrix_det}
若矩阵是方阵，则给出其行列式；否则为 \code{nil}。矩阵的系数可以是复数。

\vspace{.5em}
\begin{minipage}{.6\textwidth}
\begin{verbatim}
\directlua{
  init_elements()
  M.s = matrix.square(3, 1, 1, 0, 2, -1, -2, 1, -1, 2)
  M.s:print()
  tex.print ('\\\\')
  tex.print ("Its determinant is:  " .. M.s.det)
  }
\end{verbatim}
\end{minipage}
\begin{minipage}{.4\textwidth}
\directlua{
 init_elements()
 M.s = matrix.square(3, 1, 1, 0, 2, -1, -2, 1, -1, 2)
 M.s:print()
 tex.print('\\\\')
 tex.print("Its determinant is:  "..M.s.det)
 tex.print('\\\\')
 local a = point(1, -2)
 local b = point(0, 1)
 local c = point(1, 1)
 local d = point(1, -1)
  M.A = matrix({ { a, b }, {c, d } })
 tex.print ("Its determinant is:  "..tostring(M.A.det))}
\end{minipage}

\subsection{矩阵的元方法}

矩阵必须满足一定条件才能进行某些运算。

\begin{center}
  \bgroup
  \catcode`_=12
  \small
  \captionof{table}{矩阵的元方法。}\label{matrix:metamethods}
  \begin{tabular}{ll}
    \toprule
    \textbf{元方法} & \textbf{参考} \\
    \midrule
    \tkzMeta{matrix}{add(M1,M2)} & 见 [\ref{ssub:addition_of_matrices}] \\
    \tkzMeta{matrix}{sub(M1,M2)} & 见 [\ref{ssub:addition_of_matrices}]  \\
    \tkzMeta{matrix}{unm(M}      & |- M|  \\
    \tkzMeta{matrix}{mul(M1,M2)} & [\ref{ssub:multiplication_of_matrices}]   \\
    \tkzMeta{matrix}{pow(M,n)}   & [\ref{ssub:multiplication_of_matrices}]  \\
    \tkzMeta{matrix}{tostring(M,n)} & 显示矩阵   \\
    \tkzMeta{matrix}{eq(M1,M2)}     &  真或假  \\
  \bottomrule
  \end{tabular}
  \egroup
\end{center}



\subsubsection{矩阵的加法与减法}
\label{ssub:addition_of_matrices}
为简化输入，我用了几个函数来简化显示。

\vspace{.5em}
\begin{minipage}{.6\textwidth}
\begin{verbatim}
\directlua{
  init_elements()
  M.A = matrix({ { 1, 2 }, { 2 , -1 } })
  M.B = matrix({ { -1, 0}, { 1, 3 } })
  S = M.A + M.B
  D = M.A - M.B
  dsp(M.A,'A')
  nl() nl()
  dsp(M.B,'B')
  nl() nl()
  dsp(M.S,'S') sym(" = ")
  dsp(M.A) sym(' + ') dsp(M.B)
  nl() nl()
  dsp(M.D,'D') sym(" = ")
  dsp(M.A) sym(' - ') dsp(M.B)
}
\end{verbatim}
\end{minipage}
\begin{minipage}{.4\textwidth}
\directlua{
 init_elements()

 local function dsp (M,name)
  if name then
     tex.print(name..' = ')print_matrix(M)
  else
     print_matrix(M)
  end
 end

 local  function sym(s)
   tex.print(' '..s..' ')
 end

 local function nl()
   tex.print('\\\\')
 end

  M.A = matrix({ { 1, 2 }, { 2 , -1 } })
  M.B = matrix({ { -1, 0}, { 1, 3 } })
  M.S = M.A + M.B
  M.D = M.A - M.B
  dsp(M.A,'A')
  nl() nl()
  dsp(M.B,'B')
  nl() nl()
  dsp(M.S,'S') sym(" = ") dsp(M.A) sym(' + ') dsp(M.B)
  nl() nl()
  dsp(M.D,'D') sym(" = ") dsp(M.A) sym(' - ') dsp(M.B)}
\end{minipage}

\subsubsection{矩阵的乘法与乘方}
\label{ssub:multiplication_of_matrices}
为简化输入，我用了几个函数。它们的定义可在本文档的源码部分找到。n 为大于或小于 0 的整数，或 |'T'|。

\begin{minipage}{.5\textwidth}
\begin{verbatim}
\directlua{
  init_elements()
  M.A = matrix({ { 1, 2 }, { 2 ,-1 } })
  M.B = matrix({ { -1, 0 }, { 1, 3 } })
  M.P = M.A * M.B
  M.I = M.A ^ -1
  M.C = M.A ^ 3
  M.K = 2 * M.A}
\end{verbatim}
\end{minipage}
\begin{minipage}{.5\textwidth}
\directlua{
 init_elements()

 local function dsp (M,name)
 if name then
   tex.print(name..' = ')print_matrix(M)
  else print_matrix(M) end
  end

 local  function sym (s)
    tex.print(' '..s..' ')
  end

 local function  nl  ()
    tex.print('\\\\')
 end

  M.A = matrix({ { 1, 2 }, { 2 ,-1 } })
  M.B = matrix({ { -1, 0 }, { 1, 3 } })
  M.P = M.A * M.B
  M.I = M.A ^ -1
  M.C = M.A ^ 3
  M.K = 2 * M.A
  dsp(M.P,'P') sym(" = ") dsp(M.A) sym(' * ') dsp(M.B)
  nl() nl()
  dsp(M.A ^ -1,'$A^{-1}$')
  nl() nl()
  dsp(M.K,'K')}
\end{minipage}


\subsubsection{元方法 \code{eq}}
\label{ssub:metamthod_eq}

检测两个矩阵是否相等。结果为 \textbf{true} 或 \textbf{false}。

\subsection{matrix 类的方法}

\begin{center}
  \bgroup
  \catcode`_=12
  \small
  \captionof{table}{matrix 的函数与方法。}\label{matrix:methods}
  \begin{tabular}{lll}
  \toprule
  \textbf{函数} & \texttt{参考}   & \\
  \midrule
  \tkzFct{matrix}{new(...)}
    & 见 [\ref{ssub:method_new}; \ref{sub:matrix_creation}] & \\
  \tkzFct{matrix}{matrix.square()}
    & [\ref{ssub:method_square}] & \\
  \tkzFct{matrix}{matrix.vector()}
    & [\ref{ssub:method_vector}] & \\
  \tkzFct{matrix}{matrix.row\_vector()}
    & [\ref{ssub:function_row_vector}] & \\
  \tkzFct{matrix}{matrix.create()}
    &  & \\
  \tkzFct{matrix}{matrix.identity()}
    & [\ref{ssub:method_identity}] & \\
  \tkzFct{matrix}{matrix.htm()}
    & [\ref{ssub:method_htm}] & \\
  \tkzFct{matrix}{get\_htm\_point(V)}
    & [\ref{ssub:method_get_htm_point}] & \\
  \midrule
  \textbf{方法} & \texttt{参考}   & \\
  \midrule
  \tkzMeth{matrix}{print(s,n)}
    &  & \\
  \tkzMeth{matrix}{htm\_apply(...)}
    & [\ref{ssub:method_code_htm__apply}] & \\
  \tkzMeth{matrix}{get()}
    & [\ref{ssub:get_an_element_of_a_matrix}] & \\
  \tkzMeth{matrix}{inverse()}
    & [\ref{ssub:inverse_matrix}] & \\
  \tkzMeth{matrix}{adjugate()}
    & [\ref{ssub:method_adjugate}] & \\
  \tkzMeth{matrix}{transpose()}
    & [\ref{ssub:transpose_matrix}] & \\
  \tkzMeth{matrix}{is\_diagonal()}
    & [\ref{ssub:method_is_diagonal}] & \\
  \tkzMeth{matrix}{is\_orthogonal()}
    & [\ref{ssub:method_is_orthogonal}] & \\
  \tkzMeth{matrix}{homogenization()}
    & [\ref{ssub:method_homogenization}] & \\
  \tkzMeth{matrix}{gauss\_jordan()}
    & [\ref{ssub:method_matrix_gauss_jordan}] & \\
  \tkzMeth{matrix}{rank()}
    & [\ref{ssub:method_matrix_rank}] & \\
  \tkzMeth{matrix}{augment\_right(B)}
    & [\ref{ssub:matrix-augmentation-submatrix}] & \\
  \tkzMeth{matrix}{submatrix(r1,r2,c1,c2)}
    & [\ref{ssub:matrix-augmentation-submatrix}] & \\
  \bottomrule
  \end{tabular}
  \egroup
\end{center}




\subsubsection{方法 \tkzMeth{matrix}{print}}
\label{ssub:method_print}

加载 \tkzNamePack{amsmath} 宏包后即可使用这个方法。默认选择 \tkzEnv{latex}{bmatrix} 环境，但也可以从 \tkzEnv{latex}{matrix}、\tkzEnv{latex}{pmatrix}、\tkzEnv{latex}{Bmatrix}、“vmatrix”、“Vmatrix” 中选择。另一个选项用于设置小数点后的位数。全局变量 “tkz\_dc” 用于设置小数位数。示例如下：

\vspace{.5em}
\begin{verbatim}
\directlua{
  init_elements()
  M.n = matrix({ { math.sqrt(2), math.sqrt(3) }, { math.sqrt(4), math.sqrt(5) } })
  M.n:print('pmatrix')}
\end{verbatim}

\directlua{
  init_elements()
  M.n = matrix({ { math.sqrt(2), math.sqrt(3) }, { math.sqrt(4), math.sqrt(5) } })
  tkz_dc = 3
  M.n:print('pmatrix')
}


\vspace{.5em}
你也可以用 |print_array (M)| 函数把矩阵显示为简单数组。见下例。

对于方阵，可以传入一个值列表，其第一个元素是矩阵的阶数。

\vspace{.5em}
\begin{minipage}{.5\textwidth}
\begin{verbatim}
\directlua{
init_elements()
 M.s = matrix.square(2, 1, 0, 0, 2)
 M.s:print()}
  \end{verbatim}
\end{minipage}
\begin{minipage}{.5\textwidth}
\directlua{
  init_elements()
  M.s = matrix.square(2,1,0,0,2)
  M.s:print()}
\end{minipage}

\subsubsection{函数 \tkzFct{matrix}{new}}
\label{ssub:method_new}

这是创建矩阵的主要方法。下面是一个系数为复数的 2×3 矩阵示例：

\vspace{.5em}
\begin{minipage}{.5\textwidth}
\begin{verbatim}
\directlua{
 init_elements()
 a = point(1, 0)
 b = point(1, 1)
 c = point(-1, 1)
 d = point(0, 1)
 e = point(1, -1)
 f = point(0, -1)
 M.n = matrix({ { a, b, c }, { d, e, f } })
 M.n:print()}
\end{verbatim}
\end{minipage}
\begin{minipage}{.5\textwidth}
\directlua{
 init_elements()
 a = point(1, 0)
 b = point(1, 1)
 c = point(-1, 1)
 d = point(0, 1)
 e = point(1, -1)
 f = point(0, -1)
 M.n = matrix({ { a, b, c }, { d, e, f } })
 M.n:print()}
\end{minipage}

\subsubsection{函数 \tkzFct{matrix}{matrix.vector}}
\label{ssub:method_vector}

列矩阵这一特殊情形（常用于表示向量）可如下处理：

\vspace{.5em}
\begin{minipage}{.5\textwidth}
\begin{verbatim}
\directlua{
 init_elements()
 M.v = matrix.vector(1, 2, 3)
 M.v:print()}
  \end{verbatim}
\end{minipage}
\begin{minipage}{.5\textwidth}
\directlua{
 init_elements()
 M.v = matrix.vector(1, 2, 3)
 M.v:print()}
\end{minipage}

\subsubsection{函数 \tkzFct{matrix}{matrix.row\_vector}}
\label{ssub:function_row_vector}

\begin{mybox}
\code{ M.rv = matrix.row\_vector (1, 2, 3)}

m.rv = \directlua{matrix.row_vector (1, 2, 3):print()}
\end{mybox}

\subsubsection{函数 \tkzFct{matrix}{matrix.create(n,m)}}
\label{ssub:function_matrix_create_n_m}

\begin{mybox}
\code{ M.c = matrix.create (2, 3)}

M.c = \directlua{matrix.create (2, 3):print()}
\end{mybox}

\subsubsection{函数 \tkzFct{matrix}{matrix.square}(liste)}
\label{ssub:method_square}

我们在介绍矩阵时已经见过这个方法。需要先给出矩阵的阶数，然后逐行给出系数。

\begin{minipage}{.5\textwidth}
\begin{verbatim}
\directlua{
 init_elements()
 M.s = matrix.square(2, 2, 3, -5, 4)
 M.s:print()}
\end{verbatim}
\end{minipage}
\begin{minipage}{.5\textwidth}
\directlua{
 init_elements()
 M.s = matrix.square(2, 2, 3, -5, 4)
 M.s:print()}
\end{minipage}

\subsubsection{函数 \tkzFct{matrix}{matrix.identity}}
\label{ssub:method_identity}

创建 3 阶单位矩阵


\begin{minipage}{.5\textwidth}
\begin{verbatim}
\directlua{
  init_elements()
  M.Id_3 = matrix.identity(3)
  M.Id_3:print()}
\end{verbatim}
\end{minipage}
\begin{minipage}{.5\textwidth}
\directlua{
  init_elements()
  M.Id_3 = matrix.identity(3)
  M.Id_3:print()}
\end{minipage}

\subsubsection{方法 \tkzMeth{matrix}{is\_orthogonal}}
\label{ssub:method_is_orthogonal}

若 $A^TA = I$（在容差 \code{tkz.epsilon} 范围内），该方法返回 \code{true}，否则返回 \code{false}。

\begin{verbatim}
\directlua{
 init_elements()
 local cos = math.cos
 local sin = math.sin
 local pi = math.pi
 M.A = matrix({ { cos(pi / 6), -sin(pi / 6) }, { sin(pi / 6), cos(pi / 6) } })
 M.A:print()
 bool = M.A:is_orthogonal()
 tex.print("\\\\")
 if bool then
 	tex.print("The matrix is orthogonal")
 else
 	tex.print("The matrix is not orthogonal")
 end
 tex.print("\\\\")
 tex.print("Test: $M.A^T = M.A^{-1} ?$")
 M.A:transpose():print()
 tex.print("=")
 M.A:inverse():print()}
\end{verbatim}

\directlua{
 init_elements()
 local cos = math.cos
 local sin = math.sin
 local pi = math.pi
 M.A = matrix({{cos(pi / 6), -sin(pi / 6)}, {sin(pi / 6), cos(pi / 6)}})
 M.A:print()
 bool = M.A:is_orthogonal()
 tex.print("\\\\")
 if bool then
 	tex.print("The matrix is orthogonal")
 else
 	tex.print("The matrix is not orthogonal")
 end
 tex.print("\\\\")
 tex.print("Test: $M.A^T = M.A^{-1} ?$")
 M.A:transpose():print()
 tex.print("=")
 M.A:inverse():print()}

\subsubsection{方法 \tkzMeth{matrix}{is\_diagonal}}
\label{ssub:method_is_diagonal}

若矩阵是对角矩阵，该方法返回 \code{true}，否则返回 \code{false}。非对角元与 \code{tkz.epsilon} 比较。

\subsubsection{函数 \tkzFct{matrix}{print\_array}}
\label{ssub:display_a_table_or_array_function_code_print_array}

我们需要显示结果，下面看看不同的显示方式，并区分数组与矩阵。

下面 $A$ 是一个数组。它可以显示为简单数组，也可以显示为矩阵，但不能使用属性，也不能用 |A:print()|，因为 $A$ 不是 \code{matrix} 类的对象。若想把数组像矩阵一样显示，可以用函数 |print_matrix|（见下例）。

\vspace{.5em}
\begin{minipage}{.5\textwidth}
\begin{verbatim}
\directlua{
 init_elements()
 A = { { 1, 2 }, { 1, -1 } }
 tex.print("A = ")
 print_array(A)
 tex.print(" or ")
 print_matrix(A)
 M.A = matrix({ { 1, 1 }, { 0, 2 } })
 tex.print("\\\\")
 tex.print("M = ")
 M.A:print()}
\end{verbatim}
\end{minipage}
 \begin{minipage}{.5\textwidth}
\directlua{
 init_elements()
 A = { { 1, 2 }, { 1, -1 } }
 tex.print("A = ")
 print_array(A)
 tex.print(" or ")
 print_matrix(A)
 M.A = matrix({ { 1, 1 }, { 0, 2 } })
 tex.print("\\\\")
 tex.print("M = ")
 M.A:print()}
 \end{minipage}

\subsubsection{方法 \tkzMeth{matrix}{get}}
\label{ssub:get_an_element_of_a_matrix}
获取矩阵的一个元素。


\begin{minipage}{.5\textwidth}
\begin{verbatim}
\directlua{
  init_elements()
  M.n = matrix{ { 1, 2 }, { 2, -1 } }
  S = M.n:get(1, 1) + M.n:get(2, 2)
  tex.print(S)}
\end{verbatim}
\end{minipage}
\begin{minipage}{.5\textwidth}
\directlua{
  init_elements()
  M.n = matrix{ { 1, 2 }, { 2, -1 } }
  S = M.n:get(1, 1) + M.n:get(2, 2)
  tex.print(S)}
\end{minipage}

\subsubsection{方法 \tkzMeth{matrix}{inverse}}
\label{ssub:inverse_matrix}

若矩阵不可逆，则报错。

\begin{minipage}{.6\textwidth}
\begin{verbatim}
\directlua{
 init_elements()
 M.A = matrix({ { 1, 2 }, { 2, -1 } })
 tex.print("Inverse of $A = $")
 M.B = M.A:inverse()
 M.B:print()}
\end{verbatim}
\end{minipage}
\begin{minipage}{.4\textwidth}
\directlua{
 init_elements()
 M.A = matrix({ { 1, 2 }, { 2, -1 } })
 tex.print("Inverse of $A = $")
 M.B = M.A:inverse()
 M.B:print()}
\end{minipage}

\subsubsection{用乘方语法求逆矩阵}
\label{ssub:inverse_matrix_with_power_syntax}

\begin{minipage}{.6\textwidth}
\begin{verbatim}
\directlua{
  init_elements()
  M.n = matrix({ { 1, 0, 1 }, { 1, 2, 1 }, { 0, -1, 2 } })
  tex.print("$M = $")  print_matrix (M.n)
  tex.print('\\\\')
  tex.print("Inverse of $M = M^{-1}$")
  tex.print('\\\\','=') print_matrix(M.n ^ -1)}
\end{verbatim}
\end{minipage}
\begin{minipage}{.4\textwidth}
\directlua{
  init_elements()
  M.n = matrix({ { 1, 0, 1 }, { 1, 2, 1 }, { 0, -1, 2 } })
  tex.print("$M = $")  print_matrix (M.n)
  tex.print('\\\\')
  tex.print("$M = M^{-1} = $")
  print_matrix(M.n ^ -1)}
\end{minipage}

\subsubsection{方法 \tkzMeth{matrix}{transpose}}
\label{ssub:transpose_matrix}

转置矩阵可用 |A: transpose ()| 或 |A^{'T'}| 访问。

\vspace{.5em}
\begin{minipage}{.6\textwidth}
\begin{verbatim}
\directlua{
 init_elements()
 M.A = matrix({ { 1, 2 }, { 2, -1 } })
 M.AT = M.A:transpose()
 tex.print("$A^{'T'} = $")
 M.AT:print()}
\end{verbatim}
\end{minipage}
\begin{minipage}{.4\textwidth}
\directlua{
 init_elements()
 M.A = matrix({ { 1, 2 }, { 2, -1 } })
 M.AT = M.A:transpose()
 tex.print("$A^{'T'} = $")
 M.AT:print()}
\end{minipage}

\vspace{.5em}
注：|(A ^'T')^'T' = A|

\subsubsection{方法 \tkzMeth{matrix}{adjugate}}
\label{ssub:method_adjugate}

\begin{minipage}{.6\textwidth}
\begin{verbatim}
\directlua{
  init_elements()
  M.N =  matrix({ {1, 0, 3}, {2, 1, 0},
                  {-1, 2, 0} })
  tex.print('N = ') print_matrix(M.N)
  tex.print('\\\\')
  M.N.a = M.N:adjugate()
  M.N.i = M.N * M.N.a
  tex.print('adj(M) = ') M.N.a:print()
  tex.print('\\\\')
  tex.print('N $\\times$ adj(N) = ')
  print_matrix(M.N.i)
  tex.print('\\\\')
  tex.print('det(N) = ')
  tex.print(M.N.det)}
\end{verbatim}
\end{minipage}
\begin{minipage}{.4\textwidth}
\directlua{
  init_elements()
  M.N =  matrix({ {1, 0, 3}, {2, 1, 0}, {-1, 2, 0} })
  tex.print('N = ') print_matrix(M.N)
  tex.print('\\\\')
  M.N.a = M.N:adjugate()
  M.N.i = M.N * M.N.a
  tex.print('adj(M) = ') M.N.a:print()
  tex.print('\\\\')
  tex.print('N $\\times$ adj(N) = ') print_matrix(M.N.i)
  tex.print('\\\\')
  tex.print('det(N) = ') tex.print(M.N.det)}
\end{minipage}
\newpage


\subsubsection{方法 \tkzMeth{matrix}{diagonalize}}
\label{ssub:diagonalization}

这个方法仅适用于 2 阶矩阵。它返回两个矩阵 $D$ 和 $P$，使得 $A = PDP^{-1}$：特征值在 $D$ 的对角线上，特征向量是 $P$ 的列。复特征值以复数形式给出。具有二重特征值且尚未对角化的矩阵会报错。

\begin{minipage}{.5\textwidth}
\begin{verbatim}
\directlua{
 init_elements()
 M.A = matrix({ { 5, -3 }, { 6, -4 } })
 tex.print("A = ")
 M.A:print()
 M.D, M.P = M.A:diagonalize()
 tex.print("D = ")
 M.D:print()
 tex.print("P = ")
 M.P:print()
 M.R = M.P ^ -1 * M.A * M.P
 tex.print("\\\\")
 tex.print("Test: $D = P^{-1}AP = $ ")
 M.R:print()
 tex.print("\\\\")
 tex.print("Verification: $P^{-1}P = $ ")
 M.T = M.P ^ -1 * M.P
 M.T:print()}
\end{verbatim}
\end{minipage}
\begin{minipage}{.5\textwidth}
\directlua{
 init_elements()
 M.A = matrix({ { 5, -3 }, { 6, -4 } })
 tex.print("A = ")
 M.A:print()
 M.D, M.P = M.A:diagonalize()
 tex.print("D = ")
 M.D:print()
 tex.print("P = ")
 M.P:print()
 M.R = M.P ^ -1 * M.A * M.P
 tex.print("\\\\")
 tex.print("Test: $D = P^{-1}AP = $ ")
 M.R:print()
 tex.print("\\\\")
 tex.print("Verification: $P^{-1}P = $ ")
 M.T = M.P ^ -1 * M.P
 M.T:print()}
\end{minipage}

\subsubsection{方法 \tkzMeth{matrix}{homogenization}}
\label{ssub:method_homogenization}

\code{homogenization} 的目标是能够使用齐次变换矩阵。

取一个点 $A$，使 |z.A = point(2,-1)|。为应用 \code{htm} 矩阵，我们需要对这个点做一些处理。第一步是确定与该点关联的向量（矩阵）。这很直接，因为 point 有一个叫 \code{mtx} 的属性给出这个向量：

\begin{mybox}
z.A = point(2,0)\\
M.V = z.A.mtx:homogenization()
\end{mybox}
得到：

\begin{minipage}{.5\textwidth}
\begin{verbatim}
\directlua{
 init_elements()
 pi = math.pi
 M.h = matrix.htm(pi / 4, 3, 1)
 z.A = point(2, 0)
 M.V = z.A.mtx:homogenization()
 z.A.mtx:print()
 tex.print("then after homogenization: ")
 M.V:print()}
\end{verbatim}
\end{minipage}
\begin{minipage}{.5\textwidth}
\directlua{
 init_elements()
 pi = math.pi
 M.h = matrix.htm(pi / 4, 3, 1)
 z.A = point(2, 0)
 M.V = z.A.mtx:homogenization()
 z.A.mtx:print()
 tex.print("then after homogenization: ")
 M.V:print()}
\end{minipage}

\subsubsection{函数 \tkzFct{matrix}{matrix.htm}}
\label{ssub:method_htm}
齐次变换矩阵。

这个变换有几种使用方式。首先需要创建一个能把旋转与平移关联起来的矩阵。

主要方法是创建矩阵：

\begin{mybox}
  pi  = math.pi\\
  M.h   = matrix.htm(pi / 4, 3, 1)
\end{mybox}

这样创建一个 3×3 矩阵，结合了 $\pi/4$ 的旋转与 $\overrightarrow{t}=(3,1)$ 的平移。

完整形式为 \code{matrix.htm(phi, tx, ty, sx, sy)}：先按 (\code{sx}, \code{sy}) 缩放，再旋转角度 \code{phi}，再平移向量 (\code{tx}, \code{ty})。缩放因子默认为 1。

\directlua{
  init_elements()
  pi = math.pi
  M.h = matrix.htm(pi / 4, 3, 1)
  M.h:print()
}


现在可以应用矩阵 M。设 $A$ 为此处定义的点：\ref{ssub:method_homogenization}。通过齐次化得到列矩阵 $V$。


\begin{mybox}
M.W = M.A * M.V
\end{mybox}

\directlua{
  init_elements()
  pi  = math.pi
  M.h = matrix.htm(pi / 4 , 3 , 1)
  M.h:print()
  z.A = point(2,0)
  M.V = z.A.mtx:homogenization()
  M.V:print() tex.print('=')
  M.W = M.h * M.V
  M.W:print()
}

剩下的就是提取新点的坐标。

\subsubsection{函数 \tkzFct{matrix}{get\_htm\_point}}
\label{ssub:method_get_htm_point}

上一节我们得到了矩阵 $W$。现在需要得到它所定义的点。

函数 \code{get\_htm\_point} 从应用 \code{htm} 矩阵后得到的向量中提取一个点。

\begin{minipage}{.5\textwidth}
\begin{verbatim}
\directlua{
  init_elements()
  pi  = math.pi
  M.h   = matrix.htm(pi / 4 , 3 , 1)
  z.A = point(2,0)
  M.V = z.A.mtx:homogenization()
  M.W = M.h * M.V
  M.W:print()
  z.P = get_htm_point(M.W)
  tex.print("The affix of $P$ is: ")
  tex.print(tkz.display(z.P))}
\end{verbatim}
\end{minipage}
\begin{minipage}{.5\textwidth}
\directlua{
  init_elements()
  pi  = math.pi
  M.h   = matrix.htm(pi / 4 , 3 , 1)
  z.A = point(2,0)
  M.V = z.A.mtx:homogenization()
  M.W = M.h * M.V
  M.W:print()
  z.P = get_htm_point(M.W)
  tex.print("The affix of $P$ is: ")
  tex.print(tkz.display(z.P))}
\end{minipage}

\subsubsection{方法 \tkzMeth{matrix}{htm\_apply}}
\label{ssub:method_code_htm__apply}
上面的操作可以简化为直接在点 $A$ 上使用 \code{htm\_apply} 方法。

\begin{mybox}
|z.Ap = M: htm_apply (z.A)|\\
% 显示 (z.Ap)
\end{mybox}

方法 \code{htm\_apply} 可变换一个点、一组点或一个对象。
每个参数可以是点或对象（直线、圆、三角形、正方形、矩形、四边形、平行四边形、椭圆）。若矩阵定义的是相似变换，对象保持其类；否则圆或椭圆变为椭圆，正方形或矩形变为平行四边形。矩阵必须是 3×3 矩阵，最后一行为 (0, 0, 1)。

\directlua{
 init_elements()
 pi = math.pi
 M.h = matrix.htm(pi / 4, 3, 1)
 z.O = point(0, 0)
 z.I = point(1, 0)
 z.J = point(0, 1)
 z.A = point(2, 0)
 z.B = point(1, 2)
 L.AB = line(z.A, z.B)
 z.Op, z.Ip, z.Jp = M.h:htm_apply(z.O, z.I, z.J)
 L.ApBp = M.h:htm_apply(L.AB)
 z.Ap = L.ApBp.pa
 z.Bp = L.ApBp.pb
 z.K = point(2, 2)
 T.IJK = triangle(z.I, z.J, z.K)
 Tp = M.h:htm_apply(T.IJK)
 z.Kp = Tp.pc}

\begin{minipage}{.6\textwidth}
\begin{verbatim}
\directlua{
 init_elements()
 pi = math.pi
 M.h = matrix.htm(pi / 4, 3, 1)
 z.O = point(0, 0)
 z.I = point(1, 0)
 z.J = point(0, 1)
 z.A = point(2, 0)
 z.B = point(1, 2)
 L.AB = line(z.A, z.B)
 z.Op, z.Ip, z.Jp = M.h:htm_apply(z.O, z.I, z.J)
 L.ApBp = M.h:htm_apply(L.AB)
 z.Ap = L.ApBp.pa
 z.Bp = L.ApBp.pb
 z.K = point(2, 2)
 T.IJK = triangle(z.I, z.J, z.K)
 Tp = M.h:htm_apply(T.IJK)
 z.Kp = Tp.pc}
\end{verbatim}
\end{minipage}
\begin{minipage}{.4\textwidth}
\begin{tikzpicture}[gridded]
   \tkzGetNodes
   \tkzDrawPoints(O,O',A,B,A',B',K,K')
   \tkzLabelPoints(O,O',A,B,A',B',I,J,I',J',K,K')
   \tkzDrawSegments[->](O,I O,J O',I' O',J')
   \tkzDrawLines (A,B A',B')
   \tkzDrawPolygons[red](I,J,K I',J',K')
\end{tikzpicture}
\end{minipage}

\vspace{.5 em}

新的笛卡尔坐标系：

\vspace{.5 em}
\begin{minipage}{.5\textwidth}
\begin{verbatim}
\directlua{
 init_elements()
 pi = math.pi
 tp = tex.print
 nl = "\\\\"
 a = point(1, 0)
 b = point(0, 1)
 M.R = matrix.htm(pi / 5, 2, 1)
 M.R:print()
 tp(nl)
 M.v = matrix.vector(1, 2)
 M.v:print()
 M.v.h = M.v:homogenization()
 M.v.h:print()
 tp(nl)
 M.V = M.R * M.v.h
 M.V:print()
 z.N = get_htm_point(M.V)
tex.print(tkz.display(z.N))}
\end{verbatim}
\end{minipage}
\begin{minipage}{.5\textwidth}
\directlua{
 init_elements()
 pi = math.pi
 tp = tex.print
 nl = "\\\\"
 a = point(1, 0)
 b = point(0, 1)
 M.R = matrix.htm(pi / 5, 2, 1)
 M.R:print()
 tp(nl)
 M.v = matrix.vector(1, 2)
 M.v:print()
 M.v.h = M.v:homogenization()
 M.v.h:print()
 tp(nl)
 M.V = M.R * M.v.h
 M.V:print()
 z.N = get_htm_point(M.V)
 tex.print(tkz.display(z.N))}
\end{minipage}

\subsubsection{方法 \tkzMeth{matrix}{gauss\_jordan()}}
\label{ssub:method_matrix_gauss_jordan}


这个方法对矩阵应用高斯–约当消元算法，返回其行最简形（RREF）。

\medskip
高斯–约当算法通过一系列初等行变换把矩阵化为等价矩阵：
\begin{itemize}
\item 交换两行；
\item 用非零数乘某一行；
\item 把某一行的倍数加到另一行。
\end{itemize}
所得矩阵满足以下性质：
\begin{itemize}
\item 每个非零行的首个非零元都等于 $1$；
\item 每个首 $1$ 是其所在列唯一的非零元；
\item 所有零行（如果有）都在矩阵底部。
\end{itemize}

\medskip
\texttt{返回值：}

该方法返回一个新的矩阵对象，对应原矩阵的行最简形。
原矩阵不被修改。

\medskip
\texttt{数值说明：}

由于计算使用浮点运算，内部使用数值容差来检测零主元。
非常小的值（通常低于某个固定阈值）被视为零。

\medskip
\texttt{典型用途：}

\begin{itemize}
\item 求解线性方程组；
\item 计算矩阵的秩；
\item 检测向量组的线性无关性；
\item 计算方阵的逆（当逆存在时）。
\end{itemize}

\texttt{2×2 矩阵示例：}

\begin{tkzexample}[latex=.5\textwidth]
\directlua{
  M.A    = matrix({{2,1},{5,3}})
  M.Ainv = M.A:augment_right(matrix.identity(2))
              :gauss_jordan()
              :submatrix(1, 2, 3, 4)

  M.Ainv:print()
}
\end{tkzexample}

\texttt{3×3 矩阵示例：}

\begin{tkzexample}[latex=.5\textwidth]
\directlua{
  % Matrix 3x3
  M.A = matrix({
    {2, 1, 1},
    {1, 3, 2},
    {1, 0, 0}
  })

  % Inverse via augmentation + Gauss–Jordan
  M.Ainv = M.A
            :augment_right(matrix.identity(3))
            :gauss_jordan()
            :submatrix(1, 3, 4, 6)

  % Print the inverse
  M.Ainv:print()
}
\end{tkzexample}




\subsubsection{方法 \tkzMeth{matrix}{rank()}}
\label{ssub:method_matrix_rank}

这个方法返回矩阵的秩，即其行（等价地，列）所张成的向量空间的维数。

\medskip
秩通过高斯–约当消元过程计算：它等于矩阵行最简形中非零行的行数。

\medskip
\texttt{返回值：}
该方法返回一个非负整数，等于矩阵的秩。

\medskip
\texttt{说明：}
\begin{itemize}
\item 秩始终小于或等于行数与列数的较小者。
\item 若矩阵的秩等于这个较小者，则称矩阵满秩。
\item 方阵可逆当且仅当其秩达到最大。
\end{itemize}

\medskip
\texttt{数值稳定性：}
与 \tkzMeth{matrix}{gauss\_jordan()} 一样，使用数值容差来判断某行是否应视为零行。

\medskip
\noindent
\textbf{注。}
方法 \tkzMeth{matrix}{rank()} 内部依赖高斯–约当消元，但直接、便捷地给出秩，而不暴露中间的行最简形矩阵。

\texttt{示例 1：}

\begin{tkzexample}[latex=.5\textwidth]
\directlua{
  M.A = matrix({
    {2, 1, 1},
    {1, 3, 2},
    {1, 0, 0}
  })

  tex.print("rank(A) = ", M.A:rank())
}
\end{tkzexample}

\texttt{示例 2：}

\begin{tkzexample}[latex=.5\textwidth]
\directlua{
  M.B = matrix({
    {1, 2, 3},
    {2, 4, 6},
    {1, 1, 1}
  })

  tex.print("rank(B) = ", M.B:rank())
}
\end{tkzexample}

\texttt{示例 3：}

\begin{tkzexample}[latex=.5\textwidth]

\directlua{
  M.C = matrix({
    {1, 2, 3},
    {2, 4, 6},
    {1, 1, 1}
  })

  M.R = M.C:gauss_jordan()

  tex.print("rank(C) = ", M.C:rank())
}

\end{tkzexample}


\subsubsection{增广矩阵与子矩阵}
\label{ssub:matrix-augmentation-submatrix}

方阵求逆与线性方程组求解都基于同一个基本思想：
先处理一个\emph{增广矩阵}，再在高斯–约当消元后提取其中的一块。

\medskip

\paragraph{增广矩阵。}
设 $A$ 是 $m\times n$ 矩阵，$B$ 是 $m\times p$ 矩阵。
增广矩阵 $[A \mid B]$ 通过把 $B$ 的各列接在 $A$ 的右侧得到：

\[
[A \mid B] =
\begin{pmatrix}
A_{11} & \cdots & A_{1n} & B_{11} & \cdots & B_{1p} \\
\vdots &        & \vdots & \vdots &        & \vdots \\
A_{m1} & \cdots & A_{mn} & B_{m1} & \cdots & B_{mp}
\end{pmatrix}.
\]

在 \texttt{matrix} 类中，这个操作由方法 \texttt{augment\_right} 完成。

\medskip

\paragraph{示例（矩阵求逆）。}
设
\[
A =
\begin{pmatrix}
2 & 1\\
5 & 3
\end{pmatrix}.
\]
增广矩阵 $[A \mid I]$ 为

\[
\begin{pmatrix}
2 & 1 & 1 & 0\\
5 & 3 & 0 & 1
\end{pmatrix}.
\]

应用高斯–约当算法得到

\[
\begin{pmatrix}
1 & 0 & 3 & -1\\
0 & 1 & -5 & 2
\end{pmatrix}
=
[I \mid A^{-1}],
\]
由此可直接读出逆矩阵。

\medskip

\paragraph{子矩阵提取。}
方法 \texttt{submatrix(r1, r2, c1, c2)} 从矩阵中提取一个矩形块，由以下范围界定：

\begin{itemize}
  \item 行从 \texttt{r1} 到 \texttt{r2}，
  \item 列从 \texttt{c1} 到 \texttt{c2}。
\end{itemize}

所有下标从~$1$ 开始。

\medskip

\paragraph{\texttt{submatrix(1,2,3,4)} 的含义。}
在上例中，消元后的增广矩阵有：
\begin{itemize}
  \item $2$ 行，
  \item $4$ 列。
\end{itemize}

逆矩阵占据：
\begin{itemize}
  \item 第 $1$ 到 $2$ 行，
  \item 第 $3$ 到 $4$ 列。
\end{itemize}

因此提取
\[
\texttt{submatrix(1,2,3,4)}
\]
恰好返回块
\[
A^{-1} =
\begin{pmatrix}
3 & -1\\
-5 & 2
\end{pmatrix}.
\]

\medskip

\paragraph{一般规则。}
对于 $n$ 阶方阵，逆矩阵通过
\[
[A \mid I_n]
\;\xrightarrow{\text{Gauss--Jordan}}\;
[I_n \mid A^{-1}],
\]
得到，并用
\[
\texttt{submatrix(1, n, n+1, 2n)}.
\]
提取。

\medskip

\paragraph{小结。}
\begin{itemize}
  \item \texttt{augment\_right} 构造增广矩阵，
  \item \texttt{gauss\_jordan} 执行行消元，
  \item \texttt{submatrix} 提取所需的结果块。
\end{itemize}

\endinput
